\documentclass[UTF8,12pt]{ctexart}
\usepackage[a4paper,margin=1in]{geometry}
\usepackage{amsmath,amssymb,bm}
\usepackage{booktabs}
\usepackage{graphicx}
\usepackage{float}
\usepackage{hyperref}
\usepackage{caption}
\usepackage{enumitem}
\hypersetup{colorlinks=true,linkcolor=black,citecolor=black,urlcolor=black}
\setlist[itemize]{leftmargin=2em,itemsep=0.2em,topsep=0.2em}

\newcommand{\InputTable}[2]{%
\IfFileExists{#1}{\input{#1}}{\begin{center}\fbox{\parbox{0.88\linewidth}{#2}}\end{center}}}

\newcommand{\PlotFigure}[3]{%
\IfFileExists{#1}{\begin{figure}[H]\centering\includegraphics[width=#2\linewidth]{#1}\caption{#3}\end{figure}}{\begin{center}\fbox{\parbox{0.88\linewidth}{图像将在运行 \texttt{python src/run\_all.py --compile-pdf} 后生成：#3}}\end{center}}}

\title{期中项目报告}
\author{张研疏 2200010823}
\date{}

\begin{document}
\maketitle

\section{任务目标}

本项目研究一个基于真实市场日频数据的长仓投资组合优化问题。给定 $n$ 个流动性较高的股票资产，投资者在每个调仓日选择权重向量
\[
    x=(x_1,\ldots,x_n)^\top\in\mathbb{R}^n,
\]
其中 $x_i$ 表示投资在第 $i$ 个资产上的资金比例。本文采用全额投资和不做空假设：
\[
    \mathbf{1}^\top x=1,\qquad x\ge 0.
\]
项目目标包括五个部分：第一，获取和清洗真实市场价格数据，并构造收益率矩阵；第二，用商业求解器求解二次规划和二阶锥规划基准模型；第三，从零实现 ADMM；第四，从零实现 PDHG；第五，在严格时间顺序下进行样本外回测，并比较等权、QP 和 SOCP 策略的收益、风险、换手率和交易成本敏感性。

\section{Task1: 数据访问、清洗与参数估计}

默认资产池为 30 只高流动性美股，时间区间为 2018-01-01 至 2025-12-31。
代码通过 Yahoo Finance 接口下载日频价格数据，并使用复权收盘价构造收益率。
数据可用性审计主要检查以下内容：
\begin{enumerate}
    \item 接口是否可以正常访问，即是否能成功下载非空价格数据；
    \item 是否需要 token、访问权限或额度限制。本文代码未使用额外 API token；
    \item 关键字段是否存在。本文至少需要日期索引和复权收盘价字段；
    \item 时间跨度是否足够。2018 至 2025 年的日频数据可以支持训练集、验证集、测试集划分，也可以用于滚动回测。
\end{enumerate}

脚本会对每个资产生成数据审计结果，包括资产代码、缺失比例、首个有效日期、最后有效日期以及是否保留。清洗前资产数量为30.

数据清洗流程如下：
\begin{enumerate}
    \item 将所有资产对齐到共同交易日历，使同一行对应同一个交易日；
    \item 计算每个资产的缺失比例
    \[
        m_i=
        \frac{\#\{t:p_{t,i}\text{ is missing}\}}{N};
    \]
    \item 删除缺失比例超过阈值的资产。本文默认阈值为
    \[
        \theta_{\mathrm{miss}}=5\%;
    \]
    \item 对于退市资产、长期停牌资产或历史数据过短的资产，若其有效样本不足或缺失比例过高，则直接删除；
    \item 对少量孤立缺失价格使用前向填充；若样本开头仍存在缺失值，则删除对应日期；
    \item 使用复权收盘价计算收益率；
    \item 对每个资产的日收益率进行 1\% 和 99\% 分位数截尾，以降低极端异常值对协方差估计的影响。
\end{enumerate}

清洗后资产数量为30. 设清洗后的价格矩阵为
\[
    P=(p_{t,i})\in\mathbb{R}^{N\times n},
\]
其中 $p_{t,i}$ 是第 $i$ 个资产在第 $t$ 个交易日的复权收盘价。本文使用简单收益率
\[
    r_{t,i}=\frac{p_{t+1,i}-p_{t,i}}{p_{t,i}},
\]
得到收益率矩阵
\[
    R=(r_{t,i})\in\mathbb{R}^{T\times n}.
\]

我们使用样本均值和样本协方差为
\[
    \hat\mu=\frac{1}{T}\sum_{t=1}^T r_t,
\]
\[
    \hat\Sigma=\frac{1}{T-1}\sum_{t=1}^T (r_t-\hat\mu)(r_t-\hat\mu)^\top.
\]
为减少数值问题，实际优化中使用
\[
    \hat\Sigma_\epsilon=\hat\Sigma+\epsilon I,
\]
其中 $\epsilon=10^{-6}$。这不会改变模型的凸性，并能改善近奇异协方差矩阵造成的求解不稳定。




\section{Task2: 商业求解器基准模型}

本节使用 CVXPY 调用 MOSEK 作为商业求解器。

\subsection{二次规划模型与有效前沿}

本文采用二次效用模型作为 QP 基准：
\[
    \min_{x\in\mathbb{R}^n}\; \frac{\gamma}{2}x^\top\hat\Sigma_\epsilon x-\hat\mu^\top x
\]
\[
    \text{s.t.}\quad \mathbf{1}^\top x=1,\qquad x\ge 0.
\]
其中 $\gamma>0$ 是风险厌恶参数。目标函数第一项度量组合方差，第二项鼓励提高期望收益。当 $\gamma$ 增大时，最优解更偏向低风险；当 $\gamma$ 较小时，最优解更偏向高收益资产。由于 $\hat\Sigma_\epsilon\succeq 0$，该问题是凸二次规划。

代码在不少于 20 个 $\gamma$ 取值上用CVXPY求解该模型，并记录组合期望收益、标准差、目标值、求解状态、运行时间，如表所示。由点
\[
    \left(\sqrt{x^\top\hat\Sigma_\epsilon x},\;\hat\mu^\top x\right)
\]
组成的曲线即为 QP 有效前沿。

\InputTable{../results/tables/qp_frontier_summary.tex}{QP求解结果表}
\PlotFigure{../results/figures/efficient_frontier.png}{0.72}{QP 与 SOCP 组合在风险-收益平面中的位置}
\PlotFigure{../results/figures/qp_weights.png}{0.72}{不同 $\gamma$ 下 QP 组合权重变化}

\subsection{二阶锥规划模型}

SOCP 基准模型为标准差预算下的收益最大化：
\[
    \max_{x\in\mathbb{R}^n}\;\hat\mu^\top x
\]
\[
    \text{s.t.}\quad \|B^\top x\|_2\le \sigma_{\max},\qquad \mathbf{1}^\top x=1,
    \qquad x\ge 0,
\]
其中 $\hat\Sigma_\epsilon=BB^\top$。因此
\[
    \|B^\top x\|_2=\sqrt{x^\top\hat\Sigma_\epsilon x}
\]
表示组合标准差。该模型的目标函数线性，风险约束是二阶锥约束，因此属于二阶锥规划。代码根据 QP 前沿中的风险范围设置多个 $\sigma_{\max}$，并将 SOCP 解画在同一个风险-收益平面中。如图1所示。

QP 和 SOCP 的经济含义不同：QP 将风险作为惩罚放入目标函数，SOCP 将风险作为硬约束；QP 通过 $\gamma$ 控制风险收益权衡，SOCP 通过 $\sigma_{\max}$ 规定可接受的最大风险。两者都是凸优化问题，但 SOCP 的约束结构更直接体现标准差预算。

\InputTable{../results/tables/socp_frontier_summary.tex}{SOCP 求解结果表}

\subsection{Task4.3: 现实修正：Box constraints}

本文采用的现实修正是 box constraints。它直接限制单个资产的最大配置比例，用来控制集中持仓风险。即使在不做空条件下，基准模型仍可能把权重集中到少数几只样本均值较高的股票上，从而使组合对个股波动过于敏感。为降低这种集中度，本文对基准模型加入单资产上限约束。

设上限参数为 $u>0$。修正后的模型写为
\[
    \min_{x\in\mathbb{R}^n}\; \frac{\gamma}{2}x^\top\hat\Sigma_\epsilon x-\hat\mu^\top x
\]
\[
    \text{s.t.}\quad \mathbf{1}^\top x=1,\qquad x\ge 0,\qquad x_i\le u,\ i=1,\ldots,n.
\]
本文把单资产上限设为 10\%，即 $u=0.1$。这个修正作为硬约束进入模型，而不是惩罚项，因为它代表的是一条明确的风险控制规则：任一单个资产都不允许超过预设持仓上限。

对 QP 而言，修正后的模型为
\[
    \min_{x\in\mathbb{R}^n}\; \frac{\gamma}{2}x^\top\hat\Sigma_\epsilon x-\hat\mu^\top x
\]
\[
    \text{s.t.}\quad \mathbf{1}^\top x=1,\qquad x\ge 0,\qquad x_i\le u,
    \quad i=1,\ldots,n.
\]
新增部分只是线性不等式约束，因此修正后的问题仍然是凸二次规划。

对 SOCP 而言，修正后的模型为
\[
    \max_{x\in\mathbb{R}^n}\;\hat\mu^\top x
\]
\[
    \text{s.t.}\quad \|B^\top x\|_2\le \sigma_{\max},\qquad \mathbf{1}^\top x=1,
    \qquad x\ge 0,\qquad x_i\le u,
    \quad i=1,\ldots,n.
\]
这里新增的仍然只是线性不等式约束，因此修正后的问题仍然是凸二阶锥规划。

这个修正会同时影响 QP 和 SOCP 的可行域。由于可行域被收紧，加入 box constraints 之后，前沿上的高收益端通常会有所收缩；同时，组合权重会变得更分散，最大单资产权重被直接截断。下表中的活跃持仓数统一定义为权重大于 $10^{-4}$ 的资产个数。

\InputTable{../results/tables/realistic_modification_summary.tex}{现实修正对组合权重与风险收益的影响将在运行脚本后自动生成。}
\InputTable{../results/tables/box_qp_frontier_summary.tex}{加入 box constraints 后的 QP 前沿摘要将在运行脚本后自动生成。}
\InputTable{../results/tables/box_socp_frontier_summary.tex}{加入 box constraints 后的 SOCP 前沿摘要将在运行脚本后自动生成。}
\PlotFigure{../results/figures/box_qp_weights.png}{0.72}{加入 box constraints 后，不同 $\gamma$ 下 QP 组合权重变化}


\section{Task3: ADMM 算法}

ADMM 用于从零求解同一个核心 QP：
\[
    \min_x\; f(x)=\frac{\gamma}{2}x^\top\hat\Sigma_\epsilon x-\hat\mu^\top x,
    \qquad \mathbf{1}^\top x=1,\quad x\ge 0.
\]
定义概率单纯形
\[
    \Delta=\{z\in\mathbb{R}^n:\mathbf{1}^\top z=1,z\ge 0\}.
\]
引入分裂变量 $z$ 后，问题写为
\[
    \min_{x,z}\;\frac{\gamma}{2}x^\top\hat\Sigma_\epsilon x-\hat\mu^\top x+I_\Delta(z),
    \qquad x-z=0.
\]
采用 scaled ADMM，迭代为
\[
    x^{k+1}=(\gamma\hat\Sigma_\epsilon+\rho I)^{-1}\left(\hat\mu+\rho(z^k-u^k)\right),
\]
\[
    z^{k+1}=\operatorname{Proj}_\Delta(x^{k+1}+u^k),
\]
\[
    u^{k+1}=u^k+x^{k+1}-z^{k+1}.
\]
其中 $x$ 步是线性方程组求解，$z$ 步是到单纯形的欧氏投影。单纯形投影通过排序阈值法实现，避免调用优化求解器。收敛监控使用原始残差和对偶残差：
\[
    r_{\rm prim}^k=\|x^k-z^k\|_2,
    \qquad
    r_{\rm dual}^k=\rho\|z^k-z^{k-1}\|_2.
\]
实验测试多个 $\rho$，比较最终目标值、残差、迭代次数和运行时间，并与商业求解器在同一 $\gamma$ 下的最优目标值比较。

\InputTable{../results/tables/admm_summary.tex}{ADMM 结果将在运行脚本后自动生成。}
\PlotFigure{../results/figures/admm_convergence.png}{0.72}{不同 $\rho$ 下 ADMM 残差收敛曲线}

\section{Task4: PDHG 算法}

PDHG 同样求解核心 QP。令
\[
    f(x)=\frac{\gamma}{2}x^\top\hat\Sigma_\epsilon x-\hat\mu^\top x,
\]
并定义线性算子
\[
    Kx=\begin{pmatrix}\mathbf{1}^\top x\\-x\end{pmatrix}.
\]
约束可写为
\[
    Kx\in\{1\}\times\mathbb{R}^n_-.
\]
其 saddle-point 形式为
\[
    \min_x\max_{p\in\mathbb{R},q\in\mathbb{R}^n_+}
    \left\{\frac{\gamma}{2}x^\top\hat\Sigma_\epsilon x-\hat\mu^\top x
    +p(\mathbf{1}^\top x-1)-q^\top x\right\}.
\]
设步长 $\tau>0$ 和 $\sigma>0$ 满足
\[
    \tau\sigma\|K\|_2^2<1.
\]
由于
\[
    K^\top K=\mathbf{1}\mathbf{1}^\top+I,
\]
所以 $\|K\|_2^2=n+1$。代码采用 $\sigma=0.9/(\tau(n+1))$，保证步长条件成立。迭代公式为
\[
    p^{k+1}=p^k+\sigma(\mathbf{1}^\top\bar{x}^k-1),
\]
\[
    q^{k+1}=\operatorname{Proj}_{\mathbb{R}^n_+}(q^k-\sigma\bar{x}^k),
\]
\[
    x^{k+1}=(I+\tau\gamma\hat\Sigma_\epsilon)^{-1}
    \left(x^k-\tau(p^{k+1}\mathbf{1}-q^{k+1})+\tau\hat\mu\right),
\]
\[
    \bar{x}^{k+1}=x^{k+1}+(x^{k+1}-x^k).
\]
实验测试多个 $\tau$，报告最终目标值、可行性残差、迭代次数和运行时间，并与 ADMM 和商业求解器比较。

\InputTable{../results/tables/pdhg_summary.tex}{PDHG 结果将在运行脚本后自动生成。}
\PlotFigure{../results/figures/pdhg_convergence.png}{0.72}{不同步长下 PDHG 可行性残差收敛曲线}

\section{Task5: 真实数据回测、敏感性分析与失败案例}


\subsection{定义与实验设计}

首先把样本按照日期分为训练集、验证集和测试集。训练集用于估计模型参数，验证集用于选择 QP 的风险厌恶参数以及 SOCP 的风险预算倍率，测试集只用于最终样本外回测。下面的表展示了各阶段的起止日期和样本数。
\InputTable{../results/tables/time_split_summary.tex}


回测采用严格样本外滚动窗口。每个调仓日仅使用该日之前的 504 个交易日估计 $\hat\mu$ 和 $\hat\Sigma_\epsilon$，然后求解组合权重，并在之后的 20 个交易日持有。也就是说，训练窗口长度为 504 个交易日，调仓频率为 20 个交易日。
比较策略包括等权组合、原始 QP、加入 box constraints 后的 QP、原始 SOCP，以及加入 box constraints 后的 SOCP。交易成本在回测里被纳入，并按换手率做线性近似扣减。

设策略日收益为 $r_t^{\rm port}$，累计财富为
\[
    W_t=\prod_{s\le t}(1+r_s^{\rm port}).
\]

Sharpe ratio 用来衡量单位风险对应多少收益，在这里采用无风险利率近似为 0 的写法：
\[
    \text{Sharpe}=\frac{R_{\rm ann}}{\sigma_{\rm ann}}.
\]
如果 Sharpe ratio 较高，表示组合在承担同样波动率时获得了更高收益；如果 Sharpe ratio 为负，则表示组合的收益表现不足以补偿它承担的波动风险。、

年化收益和年化波动率分别为
\[
    R_{\rm ann}=W_T^{252/T}-1,
    \qquad
    \sigma_{\rm ann}=\sqrt{252}\,\operatorname{std}(r_t^{\rm port}).
\]

最大回撤定义为
\[
    \min_t\left(\frac{W_t}{\max_{s\le t}W_s}-1\right).
\]

换手率定义为相邻调仓权重的 $\ell_1$ 距离。交易成本按换手率线性扣减，并测试 0、5、10 和 20 bps 的成本水平。

\subsection{参数选择}

QP 的超参数是风险厌恶系数 $\gamma$。代码在预设网格
\[
    \gamma\in\{0.5,1,2,5,10,20,50,100,200,500,1000,2000,5000\}
\]
上逐个求解，并把训练集估计出来的静态权重放到验证集收益上计算表现，最后选择验证集 Sharpe ratio 最高的 $\gamma$。这就是 QP 的 model selection 规则。

SOCP 的超参数不是直接枚举 $\sigma_{\max}$ 本身，而是先定义一个风险预算倍率 $m$，再令
\[
    \sigma_{\max}=m\cdot \sigma_{\rm eq},
\]
其中 $\sigma_{\rm eq}$ 是训练集上等权组合的标准差。因此，SOCP 的风险预算倍率不是 $\sigma_{\max}$ 本身，而是决定 $\sigma_{\max}$ 的一个相对缩放系数。代码在预设倍率集合
\[
    m\in\{0.75,0.9,1.0,1.1,1.25,1.5\}
\]
上测试，并同样选择验证集 Sharpe ratio 最高的倍率。当前结果中，QP 和 box-QP 选出的最优 $\gamma$ 都是 5000，SOCP 和 box-SOCP 选出的最优风险预算倍率都是 0.9。

\subsection{回测结果展示}




回测表格汇总 0 bps 情形下各策略的累计收益、年化收益、年化波动率、Sharpe ratio、最大回撤和平均换手率， 图6给出测试期内财富随时间的变化。
可以很明显地看出，Equal Weight 在测试集上的表现优于其他优化策略，加入 box constraints 的 SOCP 次之。然而，同时也可以发现表中Equal Weight 的年化波动率和最大回撤都显著高于其他策略，说明它的高收益伴随着更大的风险。
\InputTable{../results/tables/backtest_metrics_0bps.tex}{零交易成本回测结果将在运行脚本后自动生成。}
\PlotFigure{../results/figures/backtest_wealth.png}{0.72}{测试集累计财富曲线}

而图7说明当交易成本从 0 bps 提高到 20 bps 时，各策略年化收益如何变化。通过这些结果可以直接比较等权基准、QP、box-QP、SOCP 和 box-SOCP 的样本外表现。

\InputTable{../results/tables/cost_sensitivity.tex}{交易成本敏感性结果将在运行脚本后自动生成。}
\PlotFigure{../results/figures/cost_sensitivity.png}{0.72}{不同交易成本下的年化收益变化}
\subsection{敏感性分析}
我们检查 Model A2 对风险厌恶参数 $\gamma$ 的敏感性。对每个候选 $\gamma$，使用相同的滚动窗口回测方法分别计算验证期和测试期 Sharpe ratio。下图同时给出验证期 Sharpe、测试期 Sharpe、Equal Weight 测试期 Sharpe、验证期选出的 $\gamma$ 以及测试期事后最优 $\gamma$。

\PlotFigure{../results/figures/gamma_sensitivity_test.png}{0.68}{Model A2 对 $\gamma$ 的敏感性：验证期与测试期比较}



\subsection{失败案例}
首先由图8中可以看出，验证集和测试集中$\gamma$的最优参数不一致，分别为 $\gamma=5000$ 和 $\gamma=50$。这说明验证集的参数选择并没有很好地泛化到测试集，存在过拟合现象。
但是，无论是验证集中的最优 $\gamma = 5000$ 还是测试集事后最优的 $\gamma = 50$，在测试集上的sharpe都未超过 Equal Weight，说明该失败不是单纯由验证集参数选择造成的。

我们认为，问题主要来自均值--方差模型本身的样本外不稳定性：QP 组合依赖滚动窗口估计的 $\hat\mu$ 和 $\hat\Sigma$，而股票收益均值估计噪声较大；
当 $\hat\mu$ 对未来收益排序没有足够预测能力时，优化器会把估计噪声放大为持仓偏离。$\gamma$ 只能调节这种偏离的强弱，不能修正 $\hat\mu$ 的方向错误。
因此，本案例的失败不是简单换一个 $\gamma$ 就能避免的问题。



\section{总结}

本文研究了真实市场数据下的长仓投资组合优化问题。首先选取 30 只高流动性美股，下载并清洗 2018--2025 年日频价格数据，构造收益率矩阵，并在训练集上估计均值向量 $\hat\mu$ 和协方差矩阵 $\hat\Sigma$。为提高数值稳定性，本文对协方差矩阵加入了小的对角正则项。

在优化模型方面，本文分别建立了均值--方差 QP 模型和风险预算 SOCP 模型，并使用商业求解器得到基准解和有效前沿。同时，本文加入了单资产权重上限约束作为现实修正，以降低组合集中度。实验结果表明，box constraints 可以有效限制极端持仓，使组合更加分散。

在算法实现方面，本文从零实现了 ADMM 和 PDHG 两种一阶方法来求解核心 QP 问题。实验中，ADMM 和 PDHG 的目标值均能接近商业求解器结果，说明所实现算法在该问题上具有正确性和可用性。不同超参数会影响收敛速度和残差表现，因此一阶方法仍需要合适的步长或惩罚参数选择。

在样本外回测方面，本文采用严格时间顺序划分训练集、验证集和测试集，并使用滚动窗口重新估计参数和调仓。结果显示，在本样本中 Equal Weight 的测试期表现优于 QP、SOCP 及其 box 版本。敏感性分析进一步表明，验证集最优 $\gamma$ 与测试集事后最优 $\gamma$ 不一致，而且即使使用测试集事后最优参数，QP 仍未超过 Equal Weight。这说明失败并不只是参数没有选好，而是当前均值--方差框架在该市场样本上的样本外不稳定性。

总体来看，优化模型在训练期能够给出清晰的风险--收益权衡，但其样本外表现高度依赖 $\hat\mu$ 和 $\hat\Sigma$ 的估计质量。由于股票收益均值噪声较大，优化器可能把估计误差放大为权重偏离，导致实际表现不如更稳健的等权基准。后续改进可以考虑均值收缩、权重上限、换手惩罚、相对基准偏离约束，或在验证集不能稳定超过 Equal Weight 时直接回退到等权组合。
\end{document}

